\section{Design Implications and Limits}
\label{sec:implications}
\subsection{Constrain the transaction where evidence warrants it}
The diagnostic distinction guides control placement. A current-path effect motivates constraining verification geometry; a retained-state effect motivates preserving state formation or changing its interface; an interaction requires testing both. The RTX design applies target-phase attention constraints, whereas B570 binds a whole recipe within one context. These are distinct implementations informed by local observations, without a unique kernel diagnosis or an automatic selector.

The sequence of checks is equally important. Own-path shams validate restoration at the tested boundary; crosses reveal local dependence; complete registered-output checks qualify the candidate. B570's unchanged immediate decisions and changed continuation make the last step essential. Cost follows qualification: constraining fewer phases or accepting a wider draft can still fail the performance objective.

\subsection{Use both paid baselines for operating decisions}
The ordinary comparison answers whether the deployment change helps; the matched comparison answers whether adding a draft helps on the repaired path. RTX exposes a substantial draft-free target-path opportunity followed by a modest, uneven request increment. Both values are needed to choose an implementation. A faster qualified ordinary decoder could reduce the former, since the registered $R$ is not a demonstrated optimum.

The negative results also determine operating choices. Retain the measured RTX anchor rather than its slower qualified wider draft; reject the B570 recipe for latency against both strong baselines; avoid adding the measured M4 guard on an assumed control benefit; and address 358H target work or state interfaces before expecting draft-only optimization to reach parity. Each decision follows its own output and timer contract, rather than a pooled hardware ranking.

\subsection{Scope of the evidence}
The repair covers one Qwen3 target/draft family, finite panels and pinned runtime variants. RTX's project-unused panel consists of 24 selected MBPP inputs; unused refers to accessible project history, not training exposure. Changed geometry needs renewed qualification, as the failures demonstrate. No second-family result or preservation guarantee on arbitrary inputs is established.

Discrete agreement is weaker than numerical equivalence. The score/KV witnesses do not separate every reduction, precision, layout or scheduling effect, and shams cannot exhaust hidden backend state. Selected cuts and a four-switch family cannot establish universal minimality. Prior preservation methods were not implemented as same-contract baselines, so superiority over them remains untested.

Task medians and descriptive intervals leave session, thermal and infrastructure variation unresolved. The B570 old-N1 payload and some historical raw comparisons remain missing. Timer and historical energy domains differ; later direct-ID qualification cannot validate earlier text-based timing. The reproducible CPU checks cover included sources and exports, with original-host private-archive audits identified separately in Supplement \S S7.
